\subsection{Общий обзор}
\label{subsec:logical-overview}

Логическая структура ядра строится вокруг небольшого набора устойчивых типов и контрактов, которые используются во всех компонентах:
\begin{itemize}[compact]
	\item \textbf{протоколы} (\texttt{typing.Protocol}) задают интерфейсы \texttt{Distribution}, стратегий и вычислительных процедур, позволяя реализовывать их без жёсткой иерархии наследования;
	\item \textbf{типизированные идентификаторы} (например, \texttt{GenericCharacteristicName}, \texttt{FamilyName}) унифицируют адресацию характеристик и семейств;
	\item \textbf{тип распределения} (\texttt{DistributionType}) и связанные с ним признаки (\texttt{registry\_features}) предоставляют входные данные для ограничений применимости правил вывода характеристик;
	\item \textbf{носитель} (\texttt{Support}) описывает множество допустимых значений случайной величины и используется как в семействе (для вывода support из параметров), так и в вычислениях (для выбора корректных правил).
\end{itemize}

На рис.~\ref{fig:types} приведён обзор ключевых общесистемных типов. Указанные сущности выступают «языком сопряжения» между \texttt{families}, \texttt{distributions} и будущими \texttt{transformations}: правила и стратегии опираются на них, но не зависят от конкретных реализаций распределений.

\begin{figure}[H]
	\centering
	\includegraphics[width=.7\linewidth]{types.png}
	\caption{Общие типы и идентификаторы: виды распределений, имена характеристик и семейств, а также базовые представления носителя}
	\label{fig:types}
\end{figure}
